\bookchapter{結語：持好奇，持清醒}

恭喜汝讀至了這里。

這本書自手機硬件拆至軟件生態，自日常設置挖至開發者模式，自安全防護聊至數字遺產。若它幫汝解決矣一個實際問題——哪怕只是學會矣關掉一個不該開之權限、養成矣一次備份之習慣——那此字遂沒白寫。

但比任何具體技巧皆重要之，乃兩件事：

\begin{itemize}
  \item \textbf{持好奇}：手機每年皆於變，AI於變，摺疊屏於變，衛星通信於變。今天學之操作步驟明天或遂過時矣，但"欲知道何故"之好奇心永遠未能過時。多問一句"此權限乃幹何之""此數據存在矣何處""此功能將未能泄露隱私"，比記住一百個操作步驟皆有用。
  \item \textbf{持清醒}：手機乃工具，非主人。它設計出來就是爲了吸引汝之注意力、佔據汝之時間、引導汝之消費。用它乃汝之選擇，爲它用乃汝之失控。定期抬頭觀屏幕時間，定期欲"吾乃於用手機作事，抑於爲手機消磨"，這非矯情，這乃數字時代之生存本能。
\end{itemize}

技術於變，但人之需求沒變——我們須安全感、須連接、須爲尊重。手機只是一個載體。真正之“手機課”不於書裏，於汝每天解鎖屏幕之那一瞬間。

\begin{tcbpractice}{終一個練習}
合上這本書，拿起汝之手機，作一件事：\textbf{用30秒欲，汝之手機裏有哪三樣東西乃“丟矣將真正心疼”之？}而後，確保彼等爲備份矣。這本書遂值矣。

\end{tcbpractice}
